\chapter{Triển Khai}

\section{Môi trường phát triển}

\begin{table}[H]
\centering
\caption{Môi trường phát triển và công cụ sử dụng}
\label{tab:dev_env}
\begin{tabularx}{\textwidth}{l l X}
\toprule
\textbf{Thành phần} & \textbf{Phiên bản} & \textbf{Mục đích} \\
\midrule
Python        & 3.10.x    & Backend development \\
FastAPI       & 0.115.12  & Web framework \\
Motor/PyMongo & 3.x       & MongoDB driver \\
Paho-MQTT     & 2.x       & MQTT client Python \\
Uvicorn       & 0.30.x    & ASGI server \\
\midrule
Dart          & 3.x       & Flutter language \\
Flutter SDK   & 3.22.x    & Mobile framework \\
flutter\_map  & 6.x       & Leaflet map wrapper \\
\midrule
PlatformIO    & 6.x       & Embedded IDE \\
Arduino-ESP32 & 2.0.x     & ESP32 framework \\
\midrule
MongoDB       & 6.0       & Database \\
Mosquitto     & 2.0       & Local MQTT broker \\
Docker        & 24.x      & Containerization \\
\bottomrule
\end{tabularx}
\end{table}

\section{Triển khai firmware ESP32}

\subsection{Cấu trúc firmware}

Firmware được viết bằng Arduino C++ và quản lý bằng PlatformIO. Cấu hình môi trường trong \texttt{platformio.ini}:

\begin{lstlisting}[language=bash, caption=platformio.ini]
[env:esp32dev]
platform  = espressif32
board     = esp32dev
framework = arduino
monitor_speed = 115200
lib_deps  =
    knolleary/PubSubClient @ ^2.8
    bblanchon/ArduinoJson @ ^7.0
\end{lstlisting}

\subsection{Vòng lặp chính firmware}

\begin{lstlisting}[language=C++, caption=Vòng lặp chính main.cpp (rút gọn)]
void loop() {
    // (1) Duy trì kết nối MQTT
    if (!mqttClient.connected()) reconnectMQTT();
    mqttClient.loop();

    // (2) Đọc GPS định kỳ
    if (millis() - lastPublish >= PUBLISH_INTERVAL_MS) {
        GPSData gps = readGPS();          // AT+CGNSSINFO
        if (gps.valid) {
            String payload = buildPayload(gps);
            mqttClient.publish(MQTT_TOPIC, payload.c_str());
            blinkLED();
        }
        lastPublish = millis();
    }
}
\end{lstlisting}

\subsection{Quy trình khởi động thiết bị}

\begin{enumerate}
  \item \textbf{Khởi tạo UART:} Mở cổng Serial2 (GPIO 27/26) tốc độ 115200 để giao tiếp với SIM7000G.
  \item \textbf{Reset module:} Kéo chân RST (GPIO 25) xuống thấp 1 giây, rồi thả ra, chờ module khởi động (3--5 giây).
  \item \textbf{Bật GPS:} Gửi lệnh \texttt{AT+CGPS=1,1}, chờ phản hồi \texttt{OK}. Module mất 30--60 giây để bắt tín hiệu vệ tinh lần đầu (cold start).
  \item \textbf{Kết nối GPRS:} Cấu hình APN của nhà mạng (ví dụ: \texttt{v-internet} cho Viettel), thiết lập kết nối dữ liệu.
  \item \textbf{Kết nối MQTT/TLS:} Kết nối đến HiveMQ Cloud (cổng 8883, TLS 1.2) với thông tin xác thực.
\end{enumerate}

\subsection{Đọc và parse dữ liệu GPS}

Kết quả lệnh \texttt{AT+CGNSSINFO} trả về định dạng:
\begin{lstlisting}[caption=Ví dụ phản hồi AT+CGNSSINFO]
+CGNSSINFO: 2,09,00,00,2101.71140,N,10551.25368,E,
            100624,151523.0,69.5,0.0,0.0,1
\end{lstlisting}

\textbf{Giải thích các trường:} mode, số vệ tinh GPS/GLONASS/BeiDou, vĩ độ (DDmm.mmmm), hướng N/S, kinh độ (DDDmm.mmmm), hướng E/W, ngày (DDMMYY), giờ UTC (HHMMSS), độ cao (m), tốc độ (knot), hướng di chuyển, HDOP.

\begin{lstlisting}[language=C++, caption=Hàm parse tọa độ NMEA sang decimal degrees]
double parseNMEA(double raw, char hemi) {
    int deg = (int)(raw / 100);
    double min = raw - deg * 100.0;
    double decimal = deg + min / 60.0;
    if (hemi == 'S' || hemi == 'W') decimal = -decimal;
    return decimal;
}
// Ví dụ: 2101.71140 N -> deg=21, min=1.7114 -> 21.028523 N
\end{lstlisting}

\subsection{Định dạng payload MQTT}

\begin{lstlisting}[language=json, caption=JSON payload gửi lên MQTT]
{
  "deviceId":  "child_01",
  "lat":       21.028523,
  "lng":       105.854267,
  "timestamp": 1718000123
}
\end{lstlisting}

\section{Triển khai backend}

\subsection{Khởi động ứng dụng và vòng đời}

\begin{lstlisting}[language=python, caption=backend/app/main.py -- startup/shutdown events]
@app.on_event("startup")
async def startup():
    await mongo.connect()          # Kết nối MongoDB
    mqtt_service.start()           # Khởi động MQTT subscriber thread
    logger.info("Backend started")

@app.on_event("shutdown")
async def shutdown():
    mqtt_service.stop()
    await mongo.close()
\end{lstlisting}

\subsection{MQTT Subscriber}

\begin{lstlisting}[language=python, caption=services/mqtt\_service.py -- callback xử lý tin nhắn]
def on_message(client, userdata, msg):
    try:
        payload = json.loads(msg.payload.decode())
        asyncio.run_coroutine_threadsafe(
            location_processor.process(payload),
            asyncio.get_event_loop()
        )
    except Exception as e:
        logger.error(f"MQTT parse error: {e}")
\end{lstlisting}

\subsection{LocationProcessor -- xử lý điểm GPS}

\begin{lstlisting}[language=python, caption=Lõi xử lý trong location\_processor.py]
async def process(self, payload: dict):
    loc = LocationIn(**payload)

    # (1) Lọc nhiễu: bỏ qua nếu di chuyển < 5m so với điểm trước
    last = await location_repo.get_latest(loc.deviceId)
    if last and haversine(last, loc) < NOISE_THRESHOLD_M:
        return

    # (2) Kiểm tra geofence
    cfg = await geofence_repo.get_config()
    inside, dist, gf_id = geofence_svc.check(loc, cfg)

    # (3) Cập nhật máy trạng thái cảnh báo
    alert_event = alert_state_svc.update(loc.deviceId, inside, gf_id)

    # (4) Lưu vào MongoDB
    await location_repo.insert(loc, inside, dist, gf_id)

    # (5) Gửi thông báo nếu có cảnh báo
    if alert_event:
        await notif_svc.send(loc.deviceId, alert_event, loc)

    # (6) Broadcast qua WebSocket
    await ws_manager.broadcast_location(loc, inside, dist)
    if alert_event:
        await ws_manager.broadcast_alert(loc.deviceId, alert_event)
\end{lstlisting}

\subsection{Geofence Service -- thuật toán kiểm tra}

\begin{lstlisting}[language=python, caption=services/geofence\_service.py]
def haversine(lat1, lng1, lat2, lng2) -> float:
    R = 6_371_000  # bán kính Trái Đất (m)
    phi1, phi2 = math.radians(lat1), math.radians(lat2)
    dphi = math.radians(lat2 - lat1)
    dlam = math.radians(lng2 - lng1)
    a = (math.sin(dphi/2)**2
         + math.cos(phi1) * math.cos(phi2) * math.sin(dlam/2)**2)
    return 2 * R * math.asin(math.sqrt(a))

def check(loc: LocationIn, cfg: GeofenceConfig):
    min_dist, inside, gf_id = float('inf'), False, None
    for gf in cfg.geofences:
        d = haversine(loc.lat, loc.lng, gf.centerLat, gf.centerLng)
        if d < min_dist:
            min_dist, gf_id = d, gf.id
            inside = (d <= gf.radiusM)
    return inside, min_dist, gf_id
\end{lstlisting}

\subsection{WebSocket broadcast}

\begin{lstlisting}[language=python, caption=ws/connection\_manager.py]
class ConnectionManager:
    def __init__(self):
        self.active: list[WebSocket] = []

    async def connect(self, ws: WebSocket):
        await ws.accept()
        self.active.append(ws)

    async def broadcast_location(self, loc, inside, dist):
        msg = json.dumps({
            "type": "location_update",
            "deviceId": loc.deviceId,
            "lat": loc.lat, "lng": loc.lng,
            "timestamp": loc.timestamp,
            "insideGeofence": inside,
            "distanceM": dist
        })
        for ws in list(self.active):
            try:
                await ws.send_text(msg)
            except:
                self.active.remove(ws)
\end{lstlisting}

\subsection{Notification Service}

\begin{lstlisting}[language=python, caption=Gửi cảnh báo qua Email và Telegram]
async def send_email(device_id: str, event: str, loc: LocationIn):
    subject = f"[ChildTrack] Cảnh báo: {EVENT_LABELS[event]}"
    body = (f"Thiết bị: {device_id}\n"
            f"Sự kiện: {EVENT_LABELS[event]}\n"
            f"Vị trí: {loc.lat:.6f}, {loc.lng:.6f}\n"
            f"Thời gian: {datetime.fromtimestamp(loc.timestamp)}")
    # Gửi qua SMTP (port 587, STARTTLS)
    with smtplib.SMTP(SMTP_HOST, 587) as server:
        server.starttls()
        server.login(SMTP_USER, SMTP_PASS)
        server.sendmail(SMTP_FROM, [parent_email], msg.as_string())

async def send_telegram(device_id: str, event: str, loc: LocationIn):
    text = f"[{event}] {device_id} - {EVENT_LABELS[event]}"
    url = f"https://api.telegram.org/bot{BOT_TOKEN}/sendMessage"
    # POST request đến Telegram Bot API
\end{lstlisting}

\section{Triển khai ứng dụng Flutter}

\subsection{Cấu trúc Provider (AppProvider)}

\begin{lstlisting}[language=dart, caption=Các trường trạng thái trong AppProvider]
class AppProvider with ChangeNotifier {
  // Xác thực
  String?       _token;
  UserModel?    _currentUser;
  // Thiết bị
  List<Device>  _devices = [];
  // Vị trí real-time
  Map<String, LocationData> _locations = {};
  // Geofence
  GeofenceState? _geofenceState;
  // Cảnh báo
  List<Alert>   _alerts = [];
  // Trạng thái kết nối
  SocketStatus  _socketStatus = SocketStatus.disconnected;
  // ...getters, setters, business methods
}
\end{lstlisting}

\subsection{Hiển thị bản đồ thời gian thực}

\begin{lstlisting}[language=dart, caption=Widget bản đồ với marker và geofence circle]
FlutterMap(
  options: MapOptions(
    initialCenter: LatLng(21.0285, 105.8542),
    initialZoom: 16,
    onTap: (_, latlng) => onMapTap(latlng),  // Đặt tâm geofence
  ),
  children: [
    TileLayer(urlTemplate:
      "https://tile.openstreetmap.org/{z}/{x}/{y}.png"),

    // Vòng tròn geofence
    if (geofenceState != null)
      CircleLayer(circles: [
        CircleMarker(
          point: LatLng(geofenceState!.centerLat,
                        geofenceState!.centerLng),
          radius: geofenceState!.radiusM,
          useRadiusInMeter: true,
          color: Colors.cyan.withOpacity(0.15),
          borderColor: Colors.cyan,
          borderStrokeWidth: 2,
        )
      ]),

    // Marker vị trí thiết bị
    MarkerLayer(markers: buildMarkers(locations)),

    // Đường di chuyển
    PolylineLayer(polylines: buildPolylines(history)),
  ],
)
\end{lstlisting}

\subsection{Kết nối WebSocket}

\begin{lstlisting}[language=dart, caption=SocketService -- kết nối và xử lý message]
void connect() {
  _channel = WebSocketChannel.connect(Uri.parse(wsUrl));
  _channel!.stream.listen(
    (data) => _handleMessage(jsonDecode(data)),
    onError: (_) => _startPollingFallback(),
    onDone:  () => _reconnect(),
  );
}

void _handleMessage(Map<String, dynamic> msg) {
  switch (msg['type']) {
    case 'location_update':
      provider.updateLocation(LocationData.fromJson(msg)); break;
    case 'geofence_alert':
      provider.addAlert(Alert.fromJson(msg)); break;
    case 'geofence_state_update':
      provider.updateGeofenceState(GeofenceState.fromJson(msg)); break;
  }
}
\end{lstlisting>

\section{Hạ tầng và triển khai cloud}

\subsection{Local development}

\begin{lstlisting}[language=yaml, caption=docker-compose.yml]
services:
  mongodb:
    image: mongo:6.0
    ports: ["27017:27017"]
    volumes: ["mongodb_data:/data/db"]

  mosquitto:
    image: eclipse-mosquitto:2.0
    ports: ["1883:1883", "9001:9001"]
    volumes: ["./infra/mosquitto.conf:/mosquitto/config/mosquitto.conf"]
\end{lstlisting}

\subsection{Triển khai production (Render.com + Cloud)}

\begin{figure}[H]
\centering
\begin{tikzpicture}[
  svc/.style={rectangle, draw, rounded corners=4pt, minimum width=3.5cm,
              minimum height=0.9cm, text centered, font=\small, fill=gray!8},
  arrow/.style={-Stealth, thick, draw=black!60}
]
  \node[svc, fill=orange!15] (esp) {ESP32 + SIM7000G};
  \node[svc, fill=blue!10, right=2cm of esp] (hivemq) {HiveMQ Cloud\\(MQTT, TLS 8883)};
  \node[svc, fill=green!10, right=2cm of hivemq] (render) {Render.com\\(FastAPI)};
  \node[svc, fill=purple!10, below=1cm of render] (atlas) {MongoDB Atlas\\(M0 Free)};
  \node[svc, fill=cyan!10, right=2cm of render] (flutter) {Flutter App};

  \draw[arrow] (esp)    -- node[above]{\tiny GPRS/4G} (hivemq);
  \draw[arrow] (hivemq) -- node[above]{\tiny Subscribe} (render);
  \draw[arrow] (render) -- (atlas);
  \draw[arrow] (render) -- node[above]{\tiny WS/REST} (flutter);
\end{tikzpicture}
\caption{Kiến trúc triển khai production}
\label{fig:deployment}
\end{figure}

\begin{description}
  \item[HiveMQ Cloud (Free):] MQTT broker đám mây, hỗ trợ TLS 8883, 100 kết nối đồng thời.
  \item[Render.com (Web Service):] Chạy FastAPI trên Python 3.10, auto-deploy từ GitHub.
  \item[MongoDB Atlas (M0):] 512\,MB dung lượng miễn phí, đủ cho $\sim$3 triệu điểm GPS.
\end{description}
